\documentclass[11pt]{article}
\usepackage{amsmath,amssymb,amsthm,hyperref}
\newtheorem{lemma}{Lemma}
\newcommand{\E}{\mathbb E}
\title{A bilinear Laguerre endpoint test}
\author{Research note; independently reviewed}
\date{30 September 2026}
\begin{document}\maketitle
This note isolates the algebraic improvement from the preliminary
$\Omega(n^{5/4})$ per-die bound to $\Omega(n^{4/3})$. The complete
comparison-tensor proof is in
\path{../../power_per_die_lower_bound.tex}.

\begin{lemma}[Bilinear extension]
Suppose two linear functionals $F,G$ on polynomials satisfy
\[
 |F(Q^2)-G(Q^2)|\le C\|c\|_1^2\delta
 \quad\text{whenever }Q(z)=\sum_{j=0}^d c_j L_j(\theta z).
\]
If $Q_1,Q_2$ have coefficient sums $B_1,B_2$, respectively, then
\[
 |F(Q_1Q_2)-G(Q_1Q_2)|\le2C B_1B_2\delta.
\]
\end{lemma}
\begin{proof}
The case $B_1B_2=0$ is immediate. Otherwise set
$a=\sqrt{B_2/B_1}$ and use
\[
 Q_1Q_2=\tfrac14[(aQ_1+a^{-1}Q_2)^2
                         -(aQ_1-a^{-1}Q_2)^2].
\]
Each square has coefficient sum at most $2\sqrt{B_1B_2}$.
\end{proof}

\begin{lemma}[The tapered endpoint test has small bilinear norm]
Let
\[
 T_d(y)=\sum_{j=0}^{d-1}(1-j/d)L_j(y),\qquad
 R_d(z)=\left(\frac2{\theta d}-z\right)T_d(\theta z)^2.
\]
Then $R_d=Q_dS_d$, where $Q_d=T_d(\theta z)$ has coefficient sum
$(d+1)/2$, and $S_d$ has coefficient sum at most $2/\theta$.
\end{lemma}
\begin{proof}
The Laguerre three-term recurrence gives the exact identity
\[
 yT_d(y)=\frac1d\sum_{j=0}^{d-1}L_j(y)-L_d(y).
\]
Hence $S_d=(2/(\theta d)-z)T_d(\theta z)$ has coefficients
\[
 [L_j(\theta z)]S_d=
 \begin{cases}
 (1-2j/d)/(\theta d),&0\le j<d,\\
 1/\theta,&j=d.
 \end{cases}
\]
Their absolute sum is at most $2/\theta$.
\end{proof}

The scalar endpoint comparison in the complete proof gives
$\delta=d^2/m+e^{-cm}$ uniformly for $d\le m^{1/3}$: its summation
requires only $d/\sqrt m\to0$, so the earlier provisional range
$d\le m^{1/4}$ was unnecessary. The two lemmas therefore give
\[
 |\E_\mu R_d-\beta_\ell(R_d)|
       \le C(d^3/m+d e^{-cm}).
\]
The positive scalar expectation is bounded below by an absolute
constant. Choose $d=\lfloor\kappa m^{1/3}\rfloor$ with a sufficiently
small fixed $\kappa>0$. The endpoint-forcing error is then small,
while the normalized Christoffel-square error is
$O(d^2/m)=O(1/d)$. The remaining endpoint atom argument is unchanged,
and gives $w_K,w_1=O(m^{-4/3})$, hence every equiprobable die has
$\Omega(n^{4/3})$ faces.
\end{document}
